% !TeX root = ../print.tex
\section{UC Core Model}

\newcommand{\leanid}[1]{\texttt{#1}}
\newcommand{\Mid}{\mathsf{MachineId}}
\newcommand{\EnvId}{\mathsf{env\_id}}
\newcommand{\AdvId}{\mathsf{adv\_id}}
\newcommand{\Msg}{\mathsf{Message}}
\newcommand{\Instr}{\mathsf{MessageInstruction}}
\newcommand{\Envlp}{\mathsf{Envelope}}
\newcommand{\Port}{\mathsf{CommPort}}
\newcommand{\Pay}{\mathsf{Payload}}
\newcommand{\Ctrl}{\mathsf{Controller}}
\newcommand{\Run}{\mathsf{Run}}
\newcommand{\ProbOne}{\mathsf{probTrue}}
\newcommand{\ExecDiff}{\mathsf{execDiff}}
\newcommand{\Setup}{\mathsf{ExecutionSetup}}
\newcommand{\IdealFunc}{\mathsf{IdealFunctionality}}
\newcommand{\DummyParty}{\mathsf{DummyParty}}
\newcommand{\IdealProtocol}{\mathsf{IdealProtocol}}
\newcommand{\UCEmulates}{\mathsf{UCEmulatesAt}}
\newcommand{\UCRealizes}{\mathsf{UCRealizesAt}}
\newcommand{\SecLevel}{\mathsf{SecurityLevel}}
\newcommand{\Negl}{\mathsf{Negligible}}
\newcommand{\PPTM}{\mathsf{PPT}}
\newcommand{\Corr}{\mathsf{C}}
\emergencystretch=4em

This section describes the UC core used by the current Lean development.  It is
not the EasyUC-paper interface: there are no hierarchical addresses
\(\alpha,\beta\), no separate input guard \(G\), and no interface constructor
\(\mathsf{MI}(F,A)\).  The active model is built from natural-number machine
identities, explicit communication ports, audited execution setups, and the
controller ensemble
\[
  \Ctrl.\leanid{exec}\; \sigma : \mathbb{N}\to\mathsf{PMF}\;\mathsf{Bool}.
\]
All definitions below are parametric in an arbitrary payload type \(\Pay\).

\subsection{Machines and Communication Sets}

\begin{definition}[Machine identities and ports]
Machine identities are plain natural numbers:
\[
  \Mid=\mathbb{N}, \qquad \EnvId=0,\qquad \AdvId=1 .
\]
A communication port is a triple
\[
  p=(p.\leanid{owner},p.\leanid{dest},p.\leanid{label}) : \Port
\]
whose label is one of
\[
  \leanid{input},\qquad
  \leanid{subroutineOutput},\qquad
  \leanid{backdoor}.
\]
The proof field of \(\Port\) enforces
\[
\begin{aligned}
  p.\leanid{owner} &\ne p.\leanid{dest},\\
  p.\leanid{label}=\leanid{backdoor}
    &\Longleftrightarrow
      (p.\leanid{owner}=\AdvId\ \lor\ p.\leanid{dest}=\AdvId).
\end{aligned}
\]
Thus adversarial communication is ordinary controller traffic through
backdoor-labelled ports, not a second message-delivery mechanism.
\end{definition}

\begin{definition}[Messages and envelopes]
A message carries a source identity, a label, a core instruction, and a
payload:
\[
  m=(m.\leanid{source},m.\leanid{label},
      m.\leanid{instruction},m.\leanid{payload})
    : \Msg\;\Pay .
\]
The instruction is one of
\[
  m.\leanid{instruction}\in
    \Instr=\{\leanid{plain},\ \leanid{corrupt}(pid)\}.
\]
The \(\leanid{plain}\) instruction covers ordinary protocol messages and
ordinary adversarial control messages.  The instruction
\(\leanid{corrupt}(pid)\) is an authorization-relevant request about the fixed
static corruption identity \(pid\).  It is not a port label and it does not
change which machine receives the envelope.  The core invariant is
\[
  m.\leanid{instruction}=\leanid{corrupt}(pid)
  \Longrightarrow
  m.\leanid{label}=\leanid{backdoor}.
\]
An envelope adds the destination port and proves that the port label agrees
with the message label:
\[
  e=(e.\leanid{port},e.\leanid{message}) : \Envlp\;\Pay,
  \qquad
  e.\leanid{port.label}=e.\leanid{message.label}.
\]
The destination identity is therefore read from the port, while the source
identity is stored inside the message.  This source field is not trusted by
itself; nevertheless, messages produced by honest machines normally set it to
their own identity when the protocol requires that convention.
\end{definition}

\begin{definition}[Machine programs]
A machine is
\[
  \mu=(\mu.\leanid{id},\mu.\leanid{communication\_set},
      \mu.\leanid{program}) : \leanid{Machine}\;\Pay\;O .
\]
The local program supplies an initial state sampler indexed by the security
parameter, a receive transition, a resume distribution, a halt predicate, and
an output function:
\[
\begin{array}{rcl}
\leanid{init} &:& \mathbb{N}\to \leanid{LocalState},\\
\leanid{receive} &:&
  \leanid{LocalState}\to \Msg\;\Pay\to \leanid{LocalState},\\
\leanid{resume} &:&
  \leanid{LocalState}\to
    \mathsf{PMF}\;(\leanid{ActivationResult}\;\Pay),\\
\leanid{is\_halted} &:& \leanid{LocalState}\to\mathsf{Bool},\\
\leanid{output} &:& \leanid{LocalState}\to O .
\end{array}
\]
One activation can emit at most one envelope.  A machine is well formed only if
every port in its communication set has owner \(\mu.\leanid{id}\), and the
machine has at most one port to each destination identity.  Hence the
communication set is a static, proof-carrying send capability list.
\end{definition}

\subsection{Protocols}

\begin{definition}[Protocol shape]
A protocol is a proof-carrying collection of machines:
\[
  \pi : \leanid{Protocol}\;\Pay .
\]
It contains
\[
\begin{array}{ll}
\pi.\leanid{machines} &: \hbox{the protocol machines},\\
\pi.\leanid{initial\_states} &: \mathbb{N}\to
  \mathsf{PMF}\;(\mathsf{List}\;\leanid{AnyMachineState}),\\
\pi.\leanid{corruptible\_machines} &: \hbox{the static corruption universe}.
\end{array}
\]
The field \(\pi.\leanid{corruptible\_machines}\) is only the set of machine
identities that an execution setup may choose as statically corrupted.  It is
not a global visibility set and it does not by itself grant the adversary any
extra payload semantics.  The protocol proves
\[
  \pi.\leanid{corruptible\_machines}
    \subseteq \leanid{machine\_ids}\;\pi.\leanid{machines}.
\]
The remaining fields are static obligations.  Machine identities must be
unique; \(\EnvId\) and \(\AdvId\) must not occur in the protocol; direct sends
to \(\EnvId\) are forbidden; and every protocol port whose destination is
\(\AdvId\) must be a backdoor port.
\end{definition}

\begin{definition}[Call graph relations]
The protocol file defines caller and subroutine relations from ports:
\[
\begin{array}{rcl}
\leanid{is\_caller\_of\_id}\;\mu\;j
  &\Longleftrightarrow&
  \mu \hbox{ has an } \leanid{input} \hbox{ port to } j,\\
\leanid{is\_subroutine\_of\_id}\;\mu\;j
  &\Longleftrightarrow&
  \mu \hbox{ has a } \leanid{subroutineOutput} \hbox{ port to } j.
\end{array}
\]
The fields \(\leanid{caller\_has\_matching\_subroutine}\) and
\(\leanid{subroutine\_has\_matching\_caller}\) ensure that internal protocol
calls are bidirectionally justified.  The derived predicates
\(\pi.\leanid{is\_main\_machine}\), \(\pi.\leanid{is\_internal\_machine}\), and
\(\pi.\leanid{is\_external\_identity\_of}\) are the predicates used by the
controller when it validates environment calls and routes outputs back to the
environment.
\end{definition}

\begin{definition}[Environment, adversary, and execution setup]
The environment and adversary are machines with fixed identities:
\[
  E : \leanid{Environment}\;\Pay,\qquad
  A : \leanid{Adversary}\;\Pay .
\]
The environment has identity \(\EnvId\) and a unique backdoor port to
\(\AdvId\).  The adversary has identity \(\AdvId\), a unique backdoor port to
\(\EnvId\), and no intrinsic corruption field.  A simulator is definitionally
an adversary machine:
\[
  \leanid{Simulator}\;\Pay := \leanid{Adversary}\;\Pay .
\]
For a concrete triple \((\pi,A,E)\), legal execution requires an audited setup
\[
  \sigma : \leanid{ExecutionSetup}\;\pi\;A\;E .
\]
This object replaces the old input guard and records the static corruption
pattern for this execution:
\[
  \sigma.\leanid{corrupted\_parties} : \mathsf{Finset}\;\Mid .
\]
The setup proves
\[
  \sigma.\leanid{corrupted\_parties}
    \subseteq \pi.\leanid{corruptible\_machines}.
\]
It also proves that the environment can only call main machines or its
adversary backdoor, and that the adversary's static communication-set ports
only target the environment.  Adversary-to-protocol control ports are supplied
by the runtime overlay described below, so the static adversary record does not
encode either visibility or corruption authority.
\end{definition}

\subsection{Ideal Functionalities and Ideal Protocols}

\begin{definition}[Ideal functionality]
An ideal functionality is not merely a machine.  In Lean it is a record
\[
  f : \IdealFunc\;\Pay
\]
containing the functionality machine, the party identities, and the data needed
to generate dummy parties automatically.  Its main fields are:
\[
\begin{array}{ll}
f.\leanid{party\_ids} &: \mathsf{List}\;\Mid,\\
f.\leanid{functionality\_id} &: \Mid,\\
f.\leanid{machine} &: \leanid{Machine}\;\Pay\;\mathsf{Unit},\\
f.\leanid{party\_external\_ids} &: \Mid\to\mathsf{Finset}\;\Mid,\\
f.\leanid{wrap\_input} &: \mathsf{Option}\;\Mid\to\Pay\to\Pay,\\
f.\leanid{unwrap\_output} &: \Pay\to
  \mathsf{Option}\;(\Mid\times\Pay).
\end{array}
\]
The proof fields separate parties, external identities, the functionality,
\(\EnvId\), and \(\AdvId\).  They also constrain the functionality machine's
communication set: it may send subroutine outputs to parties and backdoor
messages to the adversary, but not arbitrary traffic.
\end{definition}

\begin{definition}[Automatically generated dummy parties]
For each \(pid\in f.\leanid{party\_ids}\), Lean constructs
\[
  \leanid{mk\_dummy\_party}\;f\;pid : \DummyParty\;\Pay .
\]
The dummy party has one input port to \(f.\leanid{functionality\_id}\), one
subroutine-output port to each external identity in
\(f.\leanid{party\_external\_ids}\;pid\), and no backdoor port.  Its resume
step is fixed by the core construction:
\begin{itemize}
\item Pending \(\leanid{input}\): send
\[
  f.\leanid{wrap\_input}\;m.\leanid{source}\;m.\leanid{payload}
\]
to \(f\).
\item Pending \(\leanid{subroutineOutput}\): parse
\[
  f.\leanid{unwrap\_output}\;m.\leanid{payload}
\]
and forward the parsed payload to an external identity.
\item Pending \(\leanid{backdoor}\): drop the message.
\end{itemize}
Thus dummy parties are generated once from the ideal functionality interface;
case studies should not hand-write their forwarding behavior.  In particular,
the simulator does not directly control a dummy party.  Any information about a
corrupted party, or any command issued on behalf of that party in the ideal
world, must be represented by the ideal functionality's own program and
backdoor interface.
\end{definition}

\begin{definition}[Ideal protocol generated from a functionality]
The ideal protocol is produced by
\[
  \leanid{mk\_ideal\_protocol}\;f : \IdealProtocol\;\Pay .
\]
Its machine list is exactly
\[
  \leanid{mk\_dummy\_parties}\;f
  \quad\hbox{followed by}\quad
  f.\leanid{machine}.
\]
The generated protocol has no global visibility set.  It currently sets
\[
  \varphi_f.\leanid{corruptible\_machines}=\varnothing,
  \qquad
  \varphi_f = (\leanid{mk\_ideal\_protocol}\;f).\leanid{protocol}.
\]
The simulator can send ordinary \(\leanid{plain}\) backdoor/control messages
to the ideal functionality through the runtime overlay, and the functionality
program decides how to interpret those payloads.  A message
\(\leanid{instruction}=\leanid{corrupt}(pid)\) is also routed to the envelope
destination, not to \(pid\); the instruction only tells the controller which
static corruption identity must be authorized.

The empty \(\leanid{corruptible\_machines}\) field means the generated ideal
protocol is immediately usable for the no-corruption case
\(\Corr=\varnothing\).  Nonempty corruption patterns require an additional
corrupt-party proxy or an ideal-protocol generator that marks the appropriate
party identities as corruptible; this is a design direction, not a completed
implementation.
The constructor also proves the protocol-shape obligations.  In particular,
the generated main machines are exactly the dummy parties, and the generated
internal machine is the functionality.
\end{definition}

\subsection{Controller Execution}

\begin{definition}[Runtime communication set]
The controller does not execute directly from a machine's static
communication set.  For a setup \(\sigma\), it computes
\[
  \sigma.\leanid{runtime\_communication\_set}\;i : \mathsf{Finset}\;\Port
\]
for the currently active sender \(i\).  This set is the union of the sender's
base protocol ports and the setup-dependent overlay:
\begin{itemize}
\item \(\EnvId\) receives input ports to all main machines;
\item \(\AdvId\) receives backdoor ports to all protocol machines;
\item every protocol machine receives a backdoor port to \(\AdvId\).
\end{itemize}
This overlay is deliberately a broad control-communication overlay.  It says
which backdoor envelopes the controller can route; it does not say that a
machine is corrupted.  Ordinary \(\leanid{plain}\) backdoor messages always
pass the corruption-instruction check, and the target machine's program
decides whether such a control payload has any effect.  The project may later
add a separate control-target policy if a case study needs a narrower overlay.
\end{definition}

\begin{definition}[Outgoing-message validity]
Before routing an emitted envelope, the controller checks
\[
  \sigma.\leanid{outgoing\_message\_valid}\;i\;e .
\]
This predicate requires
\[
\begin{array}{ll}
e.\leanid{port.owner}=i, &
e.\leanid{port}\in
  \sigma.\leanid{runtime\_communication\_set}\;i .
\end{array}
\]
It also checks the source convention.  If \(i=\EnvId\), then a message sent to
\(\AdvId\) must have source \(\leanid{none}\), while a message sent to a
protocol machine must claim one of that machine's external identities.  If
\(i\ne\EnvId\), then the source must be \(\leanid{some}\;i\).  Invalid
envelopes are dropped and control returns to the environment.

The controller also checks corruption instructions:
\[
\begin{array}{rcl}
\leanid{plain} &\Longmapsto& \mathsf{True},\\
\leanid{corrupt}(pid)
  &\Longmapsto&
    e.\leanid{message.label}=\leanid{backdoor}
    \land pid\in\sigma.\leanid{corrupted\_parties}.
\end{array}
\]
The corrupted identity is read only from the instruction.  It is not inferred
from \(e.\leanid{port.owner}\), \(e.\leanid{port.dest}\), or
\(e.\leanid{message.source}\).  This is what permits the ideal-world pattern
where a simulator sends a backdoor envelope to the ideal functionality with
\(\leanid{instruction}=\leanid{corrupt}(pid)\); the functionality then decides
what leakage or commands associated with \(pid\) mean.
\end{definition}

\begin{definition}[Controller state and step]
A controller state stores the local states of \(E\), \(A\), all protocol
machines, and the currently active identity:
\[
  s=(s_E,s_A,s_\pi,i_{\mathsf{active}})
    : \leanid{ControllerState}\;\sigma .
\]
The initial active identity is always \(\EnvId\).  One controller step first
checks whether the environment has halted.  If not, it resumes the active
machine:
\[
\begin{array}{c|c}
i_{\mathsf{active}} & \hbox{step used by the controller}\\ \hline
\EnvId & \leanid{step\_environment}\\
\AdvId & \leanid{step\_adversary}\\
\hbox{protocol identity} & \leanid{step\_protocol\_machine}.
\end{array}
\]
The emitted envelope, if any, is validated and then routed to the destination
machine, the adversary, or the environment.  A subroutine output from a
non-environment sender to an external identity is delivered to the environment.
\end{definition}

\begin{definition}[Execution ensemble]
The controller runs for the fixed project budget
\(\leanid{LeanCryptoProtocols.max\_controller\_steps}\), stopping earlier if
the environment halts.  The security parameter \(n\) indexes machine
initialization and protocol initial-state sampling; it is not the step budget.
For a setup \(\sigma\), the execution ensemble is
\[
  \Run(\sigma) =
    \Ctrl.\leanid{exec}\;\sigma : \mathbb{N}\to\mathsf{PMF}\;\mathsf{Bool}.
\]
The observed bit is the environment output after the controller run:
\[
  \Pr[\Run(\sigma,n)=1]
    = \ProbOne(\Ctrl.\leanid{exec}\;\sigma\;n).
\]
For proof engineering, a path-matching proof should track the same sequence of
active identities, validates the same envelopes, and compares the same final
environment output distribution.
\end{definition}

\subsection{UC Security Definitions}

\begin{definition}[Execution difference]
For real and ideal setups over the same environment, Lean defines
\[
\begin{aligned}
  \ExecDiff(\sigma_r,\sigma_i,n)
    &=
      \left|
        \ProbOne(\Ctrl.\leanid{exec}\;\sigma_r\;n)
        -
        \ProbOne(\Ctrl.\leanid{exec}\;\sigma_i\;n)
      \right|.
\end{aligned}
\]
Perfect security is stated as equality of the full
\(\mathsf{PMF}\;\mathsf{Bool}\), not merely equality of this scalar
probability.  Statistical and computational security use \(\ExecDiff\) with a
negligible bound.
\end{definition}

\begin{definition}[Security levels]
The common security-level parameter is
\[
  \SecLevel
    = \{\leanid{perfect},\leanid{statistical},\leanid{computational}\}.
\]
The underlying indistinguishability library also defines
\(\leanid{PerfectIndist}\), \(\leanid{StatisticalIndist}\), and
\(\leanid{ComputationalIndist}\) for general ensembles.  The UC definitions
below specialize the observation to the controller output bit.
\end{definition}

\begin{definition}[UC-emulation at a static corruption pattern]
For protocols \(\pi\) and \(\varphi\) over the same payload type, Lean defines
\[
  \UCEmulates(\ell,\Corr,\pi,\varphi)
\]
by cases on the security level \(\ell\).  The finite set \(\Corr\) is the
static corruption pattern being compared.  It is a parameter of the theorem,
not a field of \(A\) or \(S\).

For fixed machines \(A,S,E\), execution is admitted only after supplying the
corresponding setup witness:
\[
  \Sigma_r(A,E)=\Setup\;\pi\;A\;E,
  \qquad
  \Sigma_i(S,E)=\Setup\;\varphi\;S\;E.
\]
These setup values are proof-carrying legality certificates for the same
machines.  The universal quantifier over \(\sigma_r\) and \(\sigma_i\) states
that the security claim is independent of the particular setup witness
supplied to Lean, as long as both setups instantiate the same static
corruption pattern:
\[
  \sigma_r.\leanid{corrupted\_parties}=\Corr,\qquad
  \sigma_i.\leanid{corrupted\_parties}=\Corr .
\]
\begin{itemize}
\item Perfect:
\[
\begin{aligned}
&\forall A\;\exists S\;\forall E,\\
&\quad
  \forall\sigma_r\!:\!\Sigma_r(A,E)\;
  \forall\sigma_i\!:\!\Sigma_i(S,E),\\
&\quad
  \sigma_r.\leanid{corrupted\_parties}=\Corr
  \Rightarrow
  \sigma_i.\leanid{corrupted\_parties}=\Corr
  \Rightarrow\\
&\quad
  \forall n,\;
  \Ctrl.\leanid{exec}\;\sigma_r\;n
  =
  \Ctrl.\leanid{exec}\;\sigma_i\;n .
\end{aligned}
\]
\item Statistical:
\[
\begin{aligned}
&\forall A\;\exists S\;\forall E,\\
&\quad
  \forall\sigma_r\!:\!\Sigma_r(A,E)\;
  \forall\sigma_i\!:\!\Sigma_i(S,E),\\
&\quad
  \sigma_r.\leanid{corrupted\_parties}=\Corr
  \Rightarrow
  \sigma_i.\leanid{corrupted\_parties}=\Corr
  \Rightarrow\\
&\quad
  \exists \epsilon,\;
  \Negl(\epsilon)\land
  \forall n,\ \ExecDiff(\sigma_r,\sigma_i,n)\le \epsilon(n).
\end{aligned}
\]
\item Computational:
\[
\begin{aligned}
&\forall A,\PPTM(A)\Rightarrow
  \exists S,\PPTM(S)\land\\
&\quad
  \forall E,\PPTM(E)\Rightarrow\\
&\quad
  \forall\sigma_r\!:\!\Sigma_r(A,E)\;
  \forall\sigma_i\!:\!\Sigma_i(S,E),\\
&\quad
  \sigma_r.\leanid{corrupted\_parties}=\Corr
  \Rightarrow
  \sigma_i.\leanid{corrupted\_parties}=\Corr
  \Rightarrow\\
&\quad
  \exists \epsilon,\;
  \Negl(\epsilon)\land
  \forall n,\ \ExecDiff(\sigma_r,\sigma_i,n)\le \epsilon(n).
\end{aligned}
\]
\end{itemize}
The no-corruption aliases
\[
\begin{array}{c}
\leanid{UCEmulatesPerfect}\\
\leanid{UCEmulatesStatistical}\\
\leanid{UCEmulatesComputational}
\end{array}
\]
instantiate \(\Corr=\varnothing\).  The corresponding
\(\leanid{WithCorruption}\) aliases expose \(\Corr\) explicitly.
\end{definition}

\begin{definition}[UC-realization at a static corruption pattern]
UC-realization is UC-emulation where the target protocol is generated from an
ideal functionality:
\[
  \UCRealizes(\ell,\Corr,\pi,f)
  \quad\hbox{uses}\quad
  \varphi_f=(\leanid{mk\_ideal\_protocol}\;f).\leanid{protocol}.
\]
For fixed \(A,S,E\), the setup types are
\[
  \Sigma_r^f(A,E)=\Setup\;\pi\;A\;E,
  \qquad
  \Sigma_i^f(S,E)=\Setup\;\varphi_f\;S\;E.
\]
There is no additional visibility premise.  If the ideal functionality needs a
simulator interface, the functionality machine must expose the relevant
backdoor/control behavior in its communication set and program.  For example,
the computational branch unfolds to
\[
\begin{aligned}
&\forall A,\PPTM(A)\Rightarrow
  \exists S,\PPTM(S)\land\\
&\quad
  \forall E,\PPTM(E)\Rightarrow\\
&\quad
  \forall\sigma_r\!:\!\Sigma_r^f(A,E)\;
  \forall\sigma_i\!:\!\Sigma_i^f(S,E),\\
&\quad
  \sigma_r.\leanid{corrupted\_parties}=\Corr
  \Rightarrow
  \sigma_i.\leanid{corrupted\_parties}=\Corr
  \Rightarrow\\
&\quad
  \exists\epsilon,\Negl(\epsilon)\land
  \forall n,\ \ExecDiff(\sigma_r,\sigma_i,n)\le\epsilon(n).
\end{aligned}
\]
The no-corruption aliases
\[
\begin{gathered}
\leanid{UCRealizesPerfect},\qquad
\leanid{UCRealizesStatistical},\\
\leanid{UCRealizesComputational}
\end{gathered}
\]
instantiate \(\Corr=\varnothing\).  The aliases ending in
\(\leanid{WithCorruption}\) keep \(\Corr\) as an explicit parameter.
\end{definition}

\subsection{Composition Theorem}

The current composition interface is expressed by
\[
  \leanid{CompositionContext}\;\Pay .
\]
A context supplies a protocol transformer \(\leanid{plug}\), projections for
the host adversary and environment, a simulator lift, and setup-level
reductions from host executions back to subprotocol executions.  The
corruption-sensitive fields are:
\[
\begin{aligned}
&(\leanid{real\_setup\_of}\;\sigma).\leanid{corrupted\_parties}
  =
  \sigma.\leanid{corrupted\_parties},\\
&(\leanid{ideal\_setup\_of}\;\sigma).\leanid{corrupted\_parties}
  =
  \sigma.\leanid{corrupted\_parties}.
\end{aligned}
\]
These equalities are the reason the theorem can preserve the same static
pattern \(\Corr\):
\[
\begin{aligned}
&\leanid{UCEmulatesAt}\;\ell\;\Corr\;\rho\;\varphi\\
&\quad\Longrightarrow
  \leanid{UCEmulatesAt}\;\ell\;\Corr\;
    (\leanid{ctx.plug}\;\rho)\;(\leanid{ctx.plug}\;\varphi).
\end{aligned}
\]
The theorem is currently an execution-reduction interface: the context must
prove exact equality between the host controller execution and the projected
subprotocol controller execution on both real and ideal sides.

\subsection{Audit}

\redcomment{TODO: to be completed.}
